\chapter{Elastic Bodies}
\label{chap:elastic-body}

\section{Why Elastic Bodies Come Before Fluids}
\label{sec:elastic-vs-fluid}

Elastic bodies and fluids share continuum balance laws, but the elastic case is
the clearer first continuum model because every material point keeps a
persistent label.  This lets us carry the configuration-space logic of the
finite-dimensional chapters almost unchanged into field theory.  A fluid has no
preferred static shear shape: its deviatoric stress is tied to the rate of
deformation.  An elastic body stores energy in deformation itself.  This chapter
therefore begins with the material map
\[
  \varphi:B\times I\to\R^3,
  \qquad
  x=\varphi(X,t),
\]
and treats \(X\in B\) as a persistent label of a material point.  The central
questions are local and material: how the deformation gradient measures stretch,
which strain fields can come from a single-valued displacement, how stored
energy produces stress, and how boundary-value problems reduce the field
equations to measurable stiffness laws.  The gauge theory of shape-changing
bodies is separated into Chapter~\ref{chap:deforming-body}, where overall rigid
motion and internal deformation become distinct variables.  Fluid mechanics is
then treated after these material-body ideas are in place, so the reader can see
exactly what changes when the persistent reference shape is no longer the main
organizing structure.

\paragraph{Roadmap.}
The chapter has one main narrative.  Deformation gradients measure local
stretching, compatibility distinguishes true displacement fields from
incompatible strain data, and a hyperelastic action gives elastodynamic
equations and stress.  The finite-strain theory is kept first because it makes
objectivity and stress measures unambiguous; the linear theory is then obtained
as a controlled small-displacement limit.

The guiding distinction is between intrinsic deformation and frame description.
Quantities such as \(C=F^TF\) and the Green strain do not change under a rigid
rotation of the observer. Stress measures such as \(P\) and \(\sigma\) live in
different configurations and are related by the deformation gradient.  The
goal is not to memorize several names for stress and strain; it is to know
which space each tensor acts on and which boundary term defines it.

\begin{assumption}[Regular deformation map]
Unless stated otherwise, \(\varphi(\cdot,t)\) is a smooth orientation-preserving
embedding of the reference body into space:
\[
  J(X,t)=\det F(X,t)>0.
\]
The inequality \(J>0\) rules out local interpenetration and orientation
reversal. Global injectivity is a stronger condition than \(J>0\); it is assumed
when contact or self-intersection would otherwise matter, but most local stress
derivations only require \(F\) to be nonsingular.
\end{assumption}

\begin{assumption}[Reference versus spatial densities]
The reference density \(\rho_0(X)\) is mass per reference volume. The spatial
density \(\rho(x,t)\) is mass per current volume. Conservation of mass for a
material volume gives
\[
  \rho(\varphi(X,t),t)J(X,t)=\rho_0(X).
\]
Thus \(J\) converts volume elements and \(\rho_0\) converts kinetic energy,
momentum balance, and stress work back to the fixed material domain \(B\).
\end{assumption}

\section{Deformation Gradient And Strain}
\label{sec:deformation-gradient-strain}

Elasticity begins by comparing nearby material points before and after
deformation. The object that does this comparison is not the displacement
itself but its derivative: two deformations that differ by a rigid translation
have the same local strain. This is why the deformation gradient \(F\), and the
metric-like tensor \(C=F^TF\), are the natural first objects.

The deformation gradient is the local linearization of the deformation map:
\[
  F^i{}_A=\frac{\partial \varphi^i}{\partial X^A}.
\]
If \(\delta X\) is a small material vector, then its spatial image is
\[
  \delta x^i=F^i{}_A\delta X^A.
\]
The right Cauchy-Green tensor measures squared lengths after deformation:
\[
  |\delta x|^2
  =
  C_{AB}\delta X^A\delta X^B,
  \qquad
  C_{AB}=F^i{}_AF^i{}_B.
\]
The Green strain is
\[
  E_{AB}=\frac12(C_{AB}-\delta_{AB}).
\]
It vanishes for every rigid motion
\[
  \varphi(X)=a+QX,\qquad Q\in\SO(3),
\]
because \(F=Q\) and \(Q^TQ=I\). This invariance is why \(E\), not \(F-I\), is a
geometric strain measure for finite deformations.

For small displacement
\[
  \varphi(X,t)=X+\xi(X,t),
\]
one has
\[
  F^i{}_A=\delta^i_A+\partial_A\xi^i.
\]
Keeping only first-order displacement gradients,
\[
  E_{AB}
  =
  \frac12(\partial_A\xi_B+\partial_B\xi_A)
  +
  O(|\nabla\xi|^2).
\]
The infinitesimal strain is therefore
\[
  \boxed{
  e_{AB}=\frac12(\partial_A\xi_B+\partial_B\xi_A).
  }
\]

\begin{figure}[htbp]
\centering
\begin{tikzpicture}[scale=1.0, >=Latex, font=\small]
  \begin{scope}
    \draw[->] (-0.2,0) -- (3.6,0) node[right] {\(X^1\)};
    \draw[->] (0,-0.2) -- (0,2.7) node[above] {\(X^2\)};
    \draw[fill=blue!7, thick]
      (0.55,0.45) .. controls (1.0,0.2) and (2.3,0.35) .. (2.85,0.75)
      .. controls (3.15,1.35) and (2.45,2.0) .. (1.65,2.15)
      .. controls (0.75,2.0) and (0.35,1.25) .. (0.55,0.45);
    \coordinate (X) at (1.4,1.15);
    \draw[black, fill=black] (X) circle (1.6pt) node[below left] {\(X\)};
    \draw[red!75!black, thick, ->] (X) -- ++(0.9,0.38)
      node[midway, above] {\(\delta X\)};
    \node at (1.75,-0.55) {reference body \(B\)};
  \end{scope}
  \draw[thick, ->] (3.8,1.35) -- node[above] {\(\varphi\)} (5.3,1.35);
  \begin{scope}[xshift=5.7cm]
    \draw[->] (-0.2,0) -- (3.7,0) node[right] {\(x^1\)};
    \draw[->] (0,-0.2) -- (0,2.7) node[above] {\(x^2\)};
    \draw[fill=green!7, thick]
      (0.55,0.55) .. controls (1.25,0.1) and (2.85,0.45) .. (3.05,1.0)
      .. controls (3.25,1.7) and (2.35,2.45) .. (1.25,2.18)
      .. controls (0.45,1.95) and (0.25,1.05) .. (0.55,0.55);
    \coordinate (x) at (1.55,1.15);
    \draw[black, fill=black] (x) circle (1.6pt) node[below left] {\(x=\varphi(X)\)};
    \draw[red!75!black, thick, ->] (x) -- ++(1.05,0.62)
      node[midway, below right, fill=white, inner sep=1pt] {\(\delta x=F\delta X\)};
    \node at (1.8,-0.55) {current placement};
  \end{scope}
\end{tikzpicture}
\caption{The deformation gradient \(F=T_X\varphi\) is the tangent map of the
deformation. It sends an infinitesimal material vector \(\delta X\) to the
spatial vector \(\delta x=F\delta X\). Strain is obtained by comparing the
lengths and angles of these vectors, not by comparing absolute positions.}
\label{fig:deformation-gradient-tangent-map}
\end{figure}

\section{Polar Decomposition And Objectivity}
\label{sec:polar-decomposition-objectivity}

The deformation gradient contains both rigid rotation and intrinsic stretch.
For \(J=\det F>0\), polar decomposition gives
\[
  \boxed{
  F=R_{\mathrm p}U=VR_{\mathrm p},
  }
\]
where \(R_{\mathrm p}\in\SO(3)\), and \(U,V\) are symmetric positive-definite
stretch tensors. The right stretch \(U\) acts on material vectors before the
polar rotation; the left stretch \(V\) acts on spatial vectors after the
rotation. The right Cauchy-Green tensor removes the rigid rotation:
\[
  C=F^TF=U^2.
\]
Similarly,
\[
  B=FF^T=V^2.
\]

\begin{derivation}[Constructing the polar factors]
Since \(F\) is nonsingular, \(C=F^TF\) is symmetric positive definite. It has a
unique symmetric positive square root
\[
  U=C^{1/2}.
\]
Define
\[
  R_{\mathrm p}=FU^{-1}.
\]
Then
\[
  R_{\mathrm p}^TR_{\mathrm p}
  =
  U^{-1}F^TFU^{-1}
  =
  U^{-1}U^2U^{-1}
  =
  I.
\]
Because \(\det F>0\) and \(\det U>0\), \(\det R_{\mathrm p}=1\). Thus
\(R_{\mathrm p}\in\SO(3)\). Finally, setting
\[
  V=R_{\mathrm p}UR_{\mathrm p}^T
\]
gives \(F=R_{\mathrm p}U=VR_{\mathrm p}\).
\end{derivation}

The eigenvalues of \(U\) are the principal stretches. They are intrinsic
measurements of local length change: a unit material vector in a principal
direction is stretched by \(\lambda_i>0\) before the rigid rotation
\(R_{\mathrm p}\) is applied. The polar rotation is not the same thing as
material spin; it is the rotation factor in the instantaneous deformation
gradient.

This distinction is the doorway to objectivity. A constitutive law for an
ordinary elastic solid may depend on stretch invariants such as those of
\(C\), but it may not change merely because an observer superposes a rigid
rotation on the whole experiment. Polar decomposition makes the separation
visible: \(R_{\mathrm p}\) tells how the local material frame is carried into
space, while \(U\) and \(C\) record the intrinsic stretching that can store
elastic energy.

\begin{figure}[htbp]
\centering
\begin{tikzpicture}[scale=1.0, >=Latex, font=\small]
  \begin{scope}
    \draw[->] (-1.25,0) -- (1.45,0) node[right] {\(E_1\)};
    \draw[->] (0,-1.15) -- (0,1.35) node[above] {\(E_2\)};
    \draw[blue!70!black, thick] (0,0) circle (0.82);
    \draw[blue!70!black, ->] (0,0) -- (0.82,0);
    \draw[blue!70!black, ->] (0,0) -- (0,0.82);
    \node[align=center] at (0,-1.55) {material unit circle};
  \end{scope}
  \draw[->, thick] (1.65,0.2) -- node[above] {\(U\)} (2.65,0.2);
  \begin{scope}[xshift=4.2cm]
    \draw[->] (-1.35,0) -- (1.55,0) node[right] {\(E_1\)};
    \draw[->] (0,-1.15) -- (0,1.35) node[above] {\(E_2\)};
    \draw[green!50!black, thick] (0,0) ellipse (1.12 and 0.55);
    \draw[green!50!black, ->] (0,0) -- (1.12,0)
      node[midway, below] {\(\lambda_1\)};
    \draw[green!50!black, ->] (0,0) -- (0,0.55)
      node[midway, left] {\(\lambda_2\)};
    \node[align=center] at (0,-1.55) {stretch ellipse};
  \end{scope}
  \draw[->, thick] (5.85,0.2) -- node[above] {\(R_{\mathrm p}\)} (6.9,0.2);
  \begin{scope}[xshift=8.6cm, rotate=28]
    \draw[green!50!black, thick] (0,0) ellipse (1.12 and 0.55);
    \draw[green!50!black, ->] (0,0) -- (1.12,0);
    \draw[green!50!black, ->] (0,0) -- (0,0.55);
  \end{scope}
  \begin{scope}[xshift=8.6cm]
    \draw[->] (-1.35,0) -- (1.55,0) node[right] {\(e_1\)};
    \draw[->] (0,-1.15) -- (0,1.35) node[above] {\(e_2\)};
    \draw[red!70!black, ->] (0.65,0.0) arc (0:28:0.65);
    \node[red!70!black] at (0.95,0.25) {\(R_{\mathrm p}\)};
    \node[align=center] at (0,-1.55) {image under \(F=R_{\mathrm p}U\)};
  \end{scope}
\end{tikzpicture}
\caption{Polar decomposition separates stretch from rigid rotation. In a
principal two-dimensional slice, \(U\) sends the material unit circle to an
ellipse with semiaxes equal to principal stretches; \(R_{\mathrm p}\) then
rotates that stretched figure without changing its lengths.}
\label{fig:polar-decomposition-stretch-rotation}
\end{figure}

This decomposition clarifies frame indifference. If an observer superposes a
rigid rotation \(Q\in\SO(3)\), the deformation gradient becomes \(QF\), but
\[
  (QF)^T(QF)=F^TF=C.
\]
Thus a frame-indifferent stored energy can be written as
\[
  W(F,X)=\widehat W(C,X).
\]
For an isotropic material, the material has no preferred direction in the
reference body, so \(\widehat W\) may depend on \(C\) only through its principal
invariants:
\[
  I_1=\tr C,\qquad
  I_2=\frac12\left((\tr C)^2-\tr(C^2)\right),\qquad
  I_3=\det C=J^2.
\]
The distinction is important. Frame indifference is invariance under rotations
of the spatial observer. Isotropy is invariance under rotations of the
material reference frame. A crystal can be frame indifferent but not isotropic.

\begin{example}[Simple shear is not a rigid rotation]
For the simple shear
\[
  \varphi(X^1,X^2,X^3)=(X^1+\gamma X^2,X^2,X^3),
\]
the deformation gradient is
\[
  F=
  \begin{pmatrix}
  1&\gamma&0\\
  0&1&0\\
  0&0&1
  \end{pmatrix}.
\]
Although \(\det F=1\), the body is not undeformed:
\[
  C=F^TF=
  \begin{pmatrix}
  1&\gamma&0\\
  \gamma&1+\gamma^2&0\\
  0&0&1
  \end{pmatrix},
  \qquad
  E_{12}=\frac{\gamma}{2}.
\]
The volume is preserved, but right angles in the reference square are changed.
This is the basic solid-mechanics distinction between volume change and shear
strain.
\end{example}

\section{Compatibility}
\label{sec:elastic-compatibility}

Not every matrix field is a deformation gradient. If
\[
  F^i{}_A=\partial_A\varphi^i
\]
for a smooth deformation, then mixed partial derivatives commute:
\[
  \boxed{
  \partial_AF^i{}_B-\partial_BF^i{}_A=0.
  }
\]
On a simply connected body, this curl-free condition is also enough to recover
\(\varphi\) from \(F\), up to a rigid translation.

For small strain, the compatibility question is subtler because \(e\) is only
the symmetric part of \(\nabla\xi\). Given a symmetric tensor field \(e_{ij}\),
we ask whether there exists a displacement \(\xi\) such that
\[
  e_{ij}=\frac12(\partial_i\xi_j+\partial_j\xi_i).
\]
The answer is controlled by the Saint-Venant tensor:
\[
  \boxed{
  \epsilon_{ik\ell}\epsilon_{jmn}
  \partial_k\partial_m e_{\ell n}=0.
  }
\]
This is often written
\[
  \operatorname{curl}\operatorname{curl} e=0.
\]
When the condition fails, the local strain field cannot be assembled into a
single-valued displacement without cuts, defects, residual stress, or a
non-Euclidean reference geometry.

\begin{remark}[Geometric meaning]
The compatibility equations are the linearized version of flatness. A strain
field defines a metric \(g_{ij}=\delta_{ij}+2e_{ij}\). If this metric is the
pullback of the Euclidean metric by an embedding, its Riemann curvature must
vanish. Linearizing that curvature gives Saint-Venant compatibility.
\end{remark}

\section{Hyperelastic Action And Piola Stress}
\label{sec:hyperelastic-action}

A hyperelastic material has stored energy density \(W(F,X)\) per reference
volume. Let \(\rho_0(X)\) be reference mass density. The action is
\[
  S[\varphi]
  =
  \int_{t_0}^{t_1}\int_B
  \left(
    \frac12\rho_0|\partial_t\varphi|^2-W(F,X)
  \right)\dd^3X\,\dd t.
\]

\begin{assumption}[Hyperelasticity]
The stress is derived from stored energy:
\[
  P^A{}_i=\frac{\partial W}{\partial F^i{}_A}.
\]
This excludes plasticity, viscosity, fracture, and history dependence. Those
effects require additional internal variables or dissipation laws.
\end{assumption}

\begin{derivation}[Elastodynamic equation]
Let \(\delta\varphi=\eta\), with \(\eta=0\) at the time endpoints. Since
\[
  \delta F^i{}_A=\partial_A\eta^i,
\]
the first variation is
\[
\begin{aligned}
  \delta S
  &=
  \int_{t_0}^{t_1}\int_B
  \left(
    \rho_0\dot\varphi_i\dot\eta_i
    -
    P^A{}_i\partial_A\eta^i
  \right)\dd^3X\,\dd t\\
  &=
  \int_{t_0}^{t_1}\int_B
  \left(
    -\rho_0\ddot\varphi_i
    +
    \partial_A P^A{}_i
  \right)\eta^i\dd^3X\,\dd t
  -
  \int_{t_0}^{t_1}\int_{\partial B}
  P^A{}_iN_A\eta^i\,\dd S\,\dd t.
\end{aligned}
\]
For arbitrary interior variations,
\[
  \boxed{
  \rho_0\ddot\varphi_i=\partial_A P^A{}_i.
  }
\]
This is the material form of momentum balance with no body force. A body force
per unit mass \(b_i\) would add \(+\rho_0b_i\) to the right-hand side. If an
external boundary traction \(t_i\) per reference area is included through
virtual work, the natural boundary condition is
\[
  P^A{}_iN_A=t_i.
\]
\end{derivation}

The derivation also explains why \(P\), not \(\sigma\), is the stress that
appears naturally in a reference-domain action: variations are differentiated
with respect to material coordinates \(X^A\), so the boundary term is a force
per reference area.

\section{Spatial Stress And Frame Indifference}
\label{sec:spatial-stress-frame-indifference}

The first Piola-Kirchhoff stress \(P\) measures force per reference area. The
Cauchy stress \(\sigma\) measures force per current spatial area. The two
stress tensors represent the same physical traction, but on different surface
elements.

\begin{derivation}[Piola transform from traction balance]
Let a material surface element have reference unit normal \(N\) and reference
area \(\dd S\). Its spatial image has unit normal \(n\) and area \(\dd a\).
Nanson's formula states
\[
  n_i\,\dd a
  =
  J(F^{-T})_{Ai}N_A\,\dd S.
\]
The force across the same cut can be written in material variables as
\[
  \dd f_i=P^A{}_iN_A\,\dd S
\]
or in spatial variables as
\[
  \dd f_i=\sigma_{ij}n_j\,\dd a.
\]
Substituting Nanson's formula gives
\[
  P^A{}_iN_A\,\dd S
  =
  \sigma_{ij}J(F^{-T})_{Aj}N_A\,\dd S.
\]
Since \(N_A\) is arbitrary,
\[
  P^A{}_i
  =
  J\sigma_{ij}(F^{-T})_{Aj}.
\]
Multiplying by \(F^k{}_A\) and using
\[
  (F^{-T})_{Aj}F^k{}_A=\delta^k{}_j
\]
gives
\[
  \boxed{
  \sigma_{ij}
  =
  \frac1J P^A{}_iF^j{}_A.
  }
\]
\end{derivation}

A stored energy is frame indifferent if a superposed rigid rotation does not
change it:
\[
  W(QF,X)=W(F,X)\qquad\text{for all }Q\in\SO(3).
\]
Let \(Q(\eps)=I+\eps A+O(\eps^2)\), with \(A^T=-A\). Differentiating
frame indifference at \(\eps=0\) gives
\[
  0=
  P^A{}_iA^i{}_jF^j{}_A.
\]
Since \(A\) is any skew matrix, the tensor \(PF^T\) must be symmetric:
\[
  P^A{}_iF^j{}_A=P^A{}_jF^i{}_A.
\]
Therefore \(\sigma=\sigma^T\). Thus symmetry of stress follows from rotational
invariance of stored energy, matching angular-momentum balance.

When \(W(F,X)=\widehat W(C,X)\), this symmetry can also be seen directly. Since
\[
  \delta C=F^T\delta F+\delta F^TF,
\]
one obtains
\[
  P=2F\frac{\partial \widehat W}{\partial C}.
\]
Because \(\partial\widehat W/\partial C\) is symmetric,
\[
  PF^T=2F\frac{\partial\widehat W}{\partial C}F^T
\]
is symmetric, and therefore so is \(\sigma=J^{-1}PF^T\).

\section{Linear Isotropic Elasticity}
\label{sec:linear-isotropic-elasticity}

For small strain in an isotropic homogeneous solid, the quadratic energy density
is
\[
  W(e)=\frac{\lambda}{2}(\tr e)^2+\mu e_{ij}e_{ij},
\]
where \(\lambda,\mu\) are Lame moduli. The stress is
\[
  \sigma_{ij}
  =
  \frac{\partial W}{\partial e_{ij}}
  =
  \lambda(\tr e)\delta_{ij}+2\mu e_{ij}.
\]
The linearized equation of motion is
\[
  \rho_0\partial_t^2\xi_i=\partial_j\sigma_{ij}+\rho_0 b_i.
\]
Substituting the stress-strain relation,
\[
\begin{aligned}
  \partial_j\sigma_{ij}
  &=
  \lambda\partial_i(\partial_k\xi_k)
  +
  \mu\partial_j(\partial_i\xi_j+\partial_j\xi_i)\\
  &=
  (\lambda+\mu)\partial_i(\nabla\cdot\xi)
  +
  \mu\Delta\xi_i.
\end{aligned}
\]
Hence
\[
  \boxed{
  \rho_0\partial_t^2\xi
  =
  (\lambda+\mu)\nabla(\nabla\cdot\xi)
  +
  \mu\Delta\xi
  +
  \rho_0 b.
  }
\]
Using
\[
  \Delta\xi=\nabla(\nabla\cdot\xi)-\nabla\times\nabla\times\xi,
\]
one sees two wave families:
\[
  c_P=\sqrt{\frac{\lambda+2\mu}{\rho_0}},
  \qquad
  c_S=\sqrt{\frac{\mu}{\rho_0}}.
\]
Longitudinal waves change volume; shear waves change shape at nearly fixed
volume.

\begin{derivation}[Plane-wave split into \(P\) and \(S\) modes]
Set \(b=0\) and try a plane wave
\[
  \xi(X,t)=a\,e^{i(k\cdot X-\omega t)},
\]
where \(a\) is a constant polarization vector and \(k\) is the wave vector.
Then
\[
  \nabla(\nabla\cdot\xi)=-k(k\cdot a)e^{i(k\cdot X-\omega t)},
  \qquad
  \Delta\xi=-|k|^2a\,e^{i(k\cdot X-\omega t)}.
\]
Substituting into the linear elastodynamic equation gives
\[
  \rho_0\omega^2 a
  =
  (\lambda+\mu)k(k\cdot a)+\mu |k|^2a.
\]
If \(a\parallel k\), then \(k(k\cdot a)=|k|^2a\), so
\[
  \omega^2=c_P^2|k|^2,
  \qquad
  c_P^2=\frac{\lambda+2\mu}{\rho_0}.
\]
If \(a\perp k\), then \(k\cdot a=0\), so
\[
  \omega^2=c_S^2|k|^2,
  \qquad
  c_S^2=\frac{\mu}{\rho_0}.
\]
Thus the eigenvectors of the acoustic tensor split into one longitudinal
polarization and two transverse polarizations in an isotropic solid.
\end{derivation}

The calculation is also a useful diagnostic for elastic constants. Stability
requires \(c_P^2>0\) and \(c_S^2>0\), while the energy argument below gives the
slightly sharper conditions \(\mu>0\) and
\(K_{\mathrm{bulk}}=\lambda+2\mu/3>0\). The two statements are consistent:
positive shear stiffness gives shear waves, and positive bulk stiffness gives a
real longitudinal wave speed.

\section{Elastic Constants And Classical Boundary-Value Problems}
\label{sec:elastic-constants-boundary-problems}

The linear isotropic equations become physically useful only after the elastic
constants and boundary conditions are tied to measurable experiments. This
section collects the standard reductions that turn the tensor equation into
one-dimensional laboratory formulas. They are not separate theories; they are
controlled approximations of the same elastic body model. This section is also
a notation checkpoint: \(E_{\mathrm Y}\), \(\nu_{\mathrm P}\), \(\lambda\),
\(\mu\), and \(K_{\mathrm{bulk}}\) are not interchangeable names for stiffness.
They measure different responses, and the boundary-value problems below show
which modulus an experiment actually sees.

\subsection{Moduli, Positivity, And Nearly Incompressible Solids}
\label{subsec:elastic-moduli-positivity}

The quadratic energy
\[
  W(e)=\frac{\lambda}{2}(\tr e)^2+\mu e:e
\]
is best understood by splitting strain into volume-changing and shape-changing
parts. Define
\[
  e_0=e-\frac13(\tr e)I,
  \qquad
  K_{\mathrm{bulk}}=\lambda+\frac23\mu .
\]
Since
\[
  e:e=e_0:e_0+\frac13(\tr e)^2,
\]
the energy becomes
\[
  \boxed{
  W(e)=\mu\,e_0:e_0+\frac12K_{\mathrm{bulk}}(\tr e)^2.
  }
\]
The shear modulus \(\mu\) penalizes distortion at fixed volume. The bulk
modulus \(K_{\mathrm{bulk}}\) penalizes volume change. For a stable isotropic
linear elastic solid one requires
\[
  \mu>0,
  \qquad
  K_{\mathrm{bulk}}>0.
\]

Engineers often use Young's modulus \(E_{\mathrm Y}\) and Poisson's ratio
\(\nu_{\mathrm P}\) instead of \(\lambda,\mu\). The conversion is
\[
  \boxed{
  E_{\mathrm Y}
  =
  \mu\frac{3\lambda+2\mu}{\lambda+\mu},
  \qquad
  \nu_{\mathrm P}
  =
  \frac{\lambda}{2(\lambda+\mu)}.
  }
\]
Equivalently,
\[
  \boxed{
  \mu=\frac{E_{\mathrm Y}}{2(1+\nu_{\mathrm P})},
  \qquad
  \lambda=
  \frac{E_{\mathrm Y}\nu_{\mathrm P}}
  {(1+\nu_{\mathrm P})(1-2\nu_{\mathrm P})},
  \qquad
  K_{\mathrm{bulk}}=
  \frac{E_{\mathrm Y}}{3(1-2\nu_{\mathrm P})}.
  }
\]
Positive energy is equivalent to
\[
  E_{\mathrm Y}>0,
  \qquad
  -1<\nu_{\mathrm P}<\frac12.
\]
Most ordinary solids have \(0<\nu_{\mathrm P}<1/2\): they contract laterally
when stretched. The limit \(\nu_{\mathrm P}\to1/2\) is the nearly
incompressible limit. In that limit \(\lambda\) and \(K_{\mathrm{bulk}}\) grow
large while \(\mu\) can remain finite, and the equations begin to resemble
incompressible fluid mechanics: pressure-like stresses enforce a volume
constraint.

\subsection{Static Navier-Cauchy Equation And Boundary Data}
\label{subsec:static-navier-cauchy}

In static linear elasticity, inertia is absent and the displacement
\(\xi(X)\) satisfies
\[
  \boxed{
  (\lambda+\mu)\nabla(\nabla\cdot\xi)+\mu\Delta\xi+\rho_0 b=0.
  }
\]
Equivalently,
\[
  \partial_j\sigma_{ij}+\rho_0b_i=0,
  \qquad
  \sigma_{ij}=\lambda(\partial_k\xi_k)\delta_{ij}
  +\mu(\partial_i\xi_j+\partial_j\xi_i).
\]
Boundary conditions are part of the problem statement. On a displacement
boundary \(\partial B_D\), one prescribes
\[
  \xi=\xi_D.
\]
On a traction boundary \(\partial B_N\), one prescribes
\[
  \sigma n=t_N,
\]
where \(n\) is the outward unit normal to the reference boundary, to linear
order. The pure traction problem determines displacement only up to rigid
translations and rotations; these rigid modes must be fixed by constraints or
by imposing zero mean displacement and zero mean rotation.

The weak form follows from virtual work. For any test displacement \(\eta\)
vanishing on \(\partial B_D\),
\[
  \int_B \sigma(\xi):e(\eta)\,\dd^3X
  =
  \int_B \rho_0 b\cdot\eta\,\dd^3X
  +
  \int_{\partial B_N} t_N\cdot\eta\,\dd S.
\]
This is the elastic analogue of deriving Cauchy's equation from balance laws:
internal stress work equals external virtual work.

\subsection{Uniaxial Extension}
\label{subsec:uniaxial-extension}

Consider a slender bar of length \(L\), cross-sectional area \(A_{\mathrm{cs}}\),
and Young modulus \(E_{\mathrm Y}\). A tensile force \(F_{\mathrm{end}}\) is
applied at the end, and lateral surfaces are traction-free. The axial stress is
\[
  \sigma_{11}=\frac{F_{\mathrm{end}}}{A_{\mathrm{cs}}},
  \qquad
  \sigma_{22}=\sigma_{33}=0.
\]
Hooke's law in the \(E_{\mathrm Y},\nu_{\mathrm P}\) form gives
\[
  e_{11}=\frac{\sigma_{11}}{E_{\mathrm Y}},
  \qquad
  e_{22}=e_{33}=-\nu_{\mathrm P}\frac{\sigma_{11}}{E_{\mathrm Y}}.
\]
Thus the axial elongation is
\[
  \boxed{
  \Delta L
  =
  \frac{F_{\mathrm{end}}L}{E_{\mathrm Y}A_{\mathrm{cs}}}.
  }
\]
The negative lateral strain is the definition of Poisson's ratio:
\[
  \nu_{\mathrm P}
  =
  -\frac{\text{lateral strain}}{\text{axial strain}}.
\]
This simple formula is often the cleanest way to remember what Young's modulus
measures: it is stress divided by axial strain in a uniaxial test with free
lateral contraction.

\subsection{Torsion Of A Circular Shaft}
\label{subsec:circular-shaft-torsion}

Let a circular shaft of radius \(a\) and length \(L\) be twisted by torque
\(T_{\mathrm{end}}\). In the Saint-Venant torsion solution, cross-sections
rotate by an angle \(\Theta(z)\) about the shaft axis. For a circular cross
section, there is no warping correction, and the twist per length is constant:
\[
  \Theta'(z)=\frac{T_{\mathrm{end}}}{\mu J_{\mathrm{polar}}},
  \qquad
  J_{\mathrm{polar}}=\frac{\pi a^4}{2}.
\]
The total twist angle is
\[
  \boxed{
  \Theta(L)-\Theta(0)=
  \frac{T_{\mathrm{end}}L}{\mu J_{\mathrm{polar}}}.
  }
\]
The shear stress grows linearly with radius:
\[
  \tau(r)=\frac{T_{\mathrm{end}}r}{J_{\mathrm{polar}}},
  \qquad
  0\leq r\leq a.
\]
The maximum shear stress occurs at the boundary \(r=a\). This is the torsional
counterpart of the axial bar formula: stiffness is the elastic modulus times a
geometric moment of the cross-section.

\subsection{Euler-Bernoulli Beam Bending}
\label{subsec:euler-bernoulli-beam}

For a slender beam, let \(x\in[0,L]\) be arclength along the undeformed neutral
axis and let \(w(x)\) be a small transverse deflection. The Euler-Bernoulli
assumptions are:
\begin{enumerate}
\item slopes are small, so curvature is \(w''(x)\) to leading order;
\item plane cross-sections remain plane and normal to the neutral axis;
\item shear deformation and rotary inertia are neglected.
\end{enumerate}
The bending energy is
\[
  U_{\mathrm{bend}}
  =
  \frac12\int_0^L E_{\mathrm Y}I_{\mathrm{area}}(w'')^2\,\dd x,
\]
where \(I_{\mathrm{area}}\) is the second moment of area of the cross-section.
With transverse distributed load \(q(x)\), the potential energy is
\[
  \Pi[w]
  =
  \frac12\int_0^L E_{\mathrm Y}I_{\mathrm{area}}(w'')^2\,\dd x
  -
  \int_0^L q(x)w(x)\,\dd x.
\]
Varying \(w\) gives
\[
\begin{aligned}
  \delta\Pi
  &=
  \int_0^L E_{\mathrm Y}I_{\mathrm{area}}w''\eta''\,\dd x
  -
  \int_0^L q\eta\,\dd x\\
  &=
  \int_0^L
  \left(E_{\mathrm Y}I_{\mathrm{area}}w''''-q\right)\eta\,\dd x
  +
  \left[
    E_{\mathrm Y}I_{\mathrm{area}}w''\eta'
    -
    E_{\mathrm Y}I_{\mathrm{area}}w'''\eta
  \right]_{0}^{L}.
\end{aligned}
\]
The beam equation is therefore
\[
  \boxed{
  E_{\mathrm Y}I_{\mathrm{area}}w''''=q.
  }
\]
The natural boundary quantities are the bending moment
\[
  M=E_{\mathrm Y}I_{\mathrm{area}}w''
\]
and shear force
\[
  V_{\mathrm{shear}}=-E_{\mathrm Y}I_{\mathrm{area}}w'''.
\]

For a cantilever clamped at \(x=0\) and loaded by a downward end force
\(P_{\mathrm{end}}\) at \(x=L\), the boundary conditions are
\[
  w(0)=0,\qquad w'(0)=0,\qquad
  M(L)=0,\qquad V_{\mathrm{shear}}(L)=P_{\mathrm{end}}.
\]
With \(q=0\) in the interior, the solution is
\[
  \boxed{
  w(x)=
  \frac{P_{\mathrm{end}}x^2(3L-x)}
  {6E_{\mathrm Y}I_{\mathrm{area}}}.
  }
\]
The tip deflection is
\[
  \boxed{
  w(L)=\frac{P_{\mathrm{end}}L^3}
  {3E_{\mathrm Y}I_{\mathrm{area}}}.
  }
\]
For a simply supported beam under uniform load \(q_0\), the deflection is
\[
  \boxed{
  w(x)=
  \frac{q_0 x(L^3-2Lx^2+x^3)}
  {24E_{\mathrm Y}I_{\mathrm{area}}}.
  }
\]
It satisfies \(w(0)=w(L)=0\), zero bending moment at both supports, and
\(E_{\mathrm Y}I_{\mathrm{area}}w''''=q_0\).

\begin{figure}[htbp]
\centering
\begin{tikzpicture}[scale=1.0, >=Latex, font=\small]
  \begin{scope}
    \draw[fill=blue!8, thick] (0,0.65) rectangle (3.5,1.15);
    \draw[thick] (0,0.55) -- (0,1.25);
    \draw[thick, ->] (3.5,0.9) -- (4.3,0.9) node[right] {\(F_{\mathrm{end}}\)};
    \draw[<->] (0,0.25) -- node[below] {\(L\)} (3.5,0.25);
    \node[align=center, fill=white, inner sep=1.5pt] at (1.75,1.72)
      {uniaxial bar\\\(\Delta L=F_{\mathrm{end}}L/(E_{\mathrm Y}A_{\mathrm{cs}})\)};
  \end{scope}
  \begin{scope}[xshift=5.7cm]
    \draw[thick] (0,1.2) -- (0,0.25);
    \foreach \y in {0.32,0.52,0.72,0.92,1.12}
      \draw[thick] (-0.18,\y-0.12) -- (0,\y);
    \draw[blue!70!black, thick, domain=0:3.6, samples=100]
      plot (\x,{1.05-0.055*\x*\x*(3.6-\x/3)});
    \draw[red!75!black, thick, ->] (3.85,0.62) -- (3.85,-0.10)
      node[below right=1pt] {\(P_{\mathrm{end}}\)};
    \draw[<->] (0,0.02) -- node[below] {\(L\)} (3.6,0.02);
    \node[align=center, fill=white, inner sep=1.5pt] at (1.85,1.72)
      {cantilever beam\\\(w(L)=P_{\mathrm{end}}L^3/(3E_{\mathrm Y}I_{\mathrm{area}})\)};
  \end{scope}
\end{tikzpicture}
\caption{Two classical reductions of three-dimensional linear elasticity. A
slender bar under axial load measures Young's modulus. A slender beam under
transverse load converts elastic energy in curvature into the fourth-order
Euler-Bernoulli equation.}
\label{fig:bar-and-beam-elasticity}
\end{figure}

\section{What To Take Forward}
\label{sec:elastic-body-take-forward}

An elastic body is described by a deformation map \(\varphi(X,t)\) from a
material reference body \(B\) into physical space. The deformation gradient
\(F^i{}_A=\partial\varphi^i/\partial X^A\) measures local stretch and rotation,
while objective strain tensors remove the arbitrary rigid rotation.
Compatibility says that not every field of strain is the strain of a
single-valued displacement; it must arise from a deformation with commuting
mixed derivatives, or equivalently satisfy the Saint-Venant conditions in the
linear theory.

The variational derivation is parallel to finite-dimensional Lagrangian
mechanics but with fields. A hyperelastic stored energy \(W(F,X)\) defines the
first Piola-Kirchhoff stress \(P^i{}_A=\partial W/\partial F^i{}_A\), and
stationarity of the action gives material momentum balance
\(\rho_0\ddot\varphi^i=\partial_A P^i{}_A+\rho_0 b^i\) together with the
correct traction boundary term. Frame indifference forces the stored energy to
depend on deformation through rotationally invariant combinations such as
\(C=F^T F\), and it implies symmetry of the spatial Cauchy stress.

Linear isotropic elasticity is the small-strain limit. Its two Lam\'e constants
encode the quadratic energy, the Navier-Cauchy equation describes displacement
waves and static equilibrium, and the classical bar, torsion, and beam problems
show how continuum equations reduce to experimentally measurable stiffness
laws. The next chapter uses this elastic language as the material side of a
gauge theory of deforming bodies.


\section{Chapter-End Laboratories And Exercises}
\label{sec:elastic-body-chapter-end-labs}
\label{sec:elastic-body-questions}

The chapter-end work is organized as an audit of the elastic modelling chain:
objectivity, variational stress, linearization, and classical reductions. The
questions are short, but each one points to a place where continuum mechanics
often becomes confusing if one forgets which configuration a tensor belongs to.

\begin{enumerate}
\item Starting from \(C=F^TF\), show explicitly that superposing a rigid spatial
rotation on \(\varphi\) leaves \(C\) unchanged. Which strain tensors would fail
this objectivity test?
\item In the hyperelastic action, identify the boundary term whose coefficient
is the first Piola-Kirchhoff traction. How does that term transform into a
spatial Cauchy traction?
\item Derive the Navier-Cauchy equation from the quadratic isotropic energy and
identify where the two Lame constants enter longitudinal and transverse waves.
\item For the cantilever beam formula, check the four boundary conditions at the
clamped and loaded ends. Which condition represents the applied point load?
\item Compare the bar and beam reductions in Figure~\ref{fig:bar-and-beam-elasticity}.
Which geometric small parameter justifies replacing a three-dimensional elastic
body by a one-dimensional model?
\end{enumerate}

\subsection{Boundary-Value Laboratory: Bar Versus Beam}
\label{sec:elastic-bar-beam-lab}

Use Figure~\ref{fig:bar-and-beam-elasticity} as a modelling checklist. The bar
problem assumes that the dominant strain is axial and approximately uniform
over each cross-section, giving the stiffness law \(F=EA\,\Delta L/L\). The
beam problem assumes that cross-sections remain approximately plane and that
curvature, not axial stretch, carries the leading elastic energy, giving
\[
  EI\,w''''(x)=q(x).
\]
Both are one-dimensional reductions of the same three-dimensional elastic
theory. The useful comparison is therefore not ``bar formula versus beam
formula,'' but which geometric small parameter justifies throwing away the
other strain components. A good solution should state the slenderness
assumption, the boundary data, the stress or moment being measured, and the
energy term that survives the reduction.

\section{Associated Demos}
\label{sec:elastic-body-demos}

The associated elastic-boundary-value demos are:
\begin{itemize}
\item \path{demos/python/linear_elasticity.py}
\item \path{demos/mathematica/LinearElasticity.wl}
\end{itemize}
They solve elementary elastic boundary-value problems and visualize
displacement, strain, and stress fields. They are designed to be read beside
this chapter's derivations of the Piola stress, Cauchy stress, and linear
elastic equations.
